% TOP_PROBLEM: {"problemId": "problem.capacity-of-the-binary-deletion-channel", "canonicalTitle": "Capacity of the Binary Deletion Channel", "releaseRank": 225, "primaryDomain": "cryptography_coding_information_optimization", "primaryDomainLabel": "Cryptography, coding, information & optimization", "displayStatus": "open", "releaseStatus": "open"}
% !TeX program = lualatex
% Definition and sources only; this file is not an LLM solution attempt.
\documentclass[11pt]{article}
\usepackage[margin=25mm]{geometry}
\usepackage{fontspec}
\setmainfont{Latin Modern Roman}
\usepackage{amsmath,amssymb,mathtools}
\usepackage{unicode-math}
\setmathfont{Latin Modern Math}
\usepackage{xurl,hyperref}
\hypersetup{colorlinks=true,urlcolor=blue}
\setcounter{secnumdepth}{0}
\setlength{\parindent}{0pt}
\setlength{\parskip}{5pt}
\setlength{\emergencystretch}{3em}
\begin{document}
\section*{\texorpdfstring{Capacity of the Binary Deletion Channel}{Capacity of the Binary Deletion Channel}}\label{225}
\subsection{Definitions and mathematical statement}
For deletion probability $0<d<1$, a binary word $x^n$ is sent through independent deletions, and the output is the surviving subsequence in its original order. Neither terminal is told the deleted positions. A code has a uniform message $M$, an $n$-bit encoding, and a decoder acting on the variable-length output $Y$. Define
\[
 C_{\mathrm{del}}(d)=\sup\{R:\exists\text{ codes with }\liminf_n n^{-1}\log_2|\mathcal M_n|\ge R,
 \ \Pr(\widehat M\ne M)\to0\}.
\]
Determine this capacity exactly for every $d$. Rates are per input bit, not per surviving bit. A finite-precision numerical approximation, or substituting erasures whose positions are known, does not answer the exact deletion-channel question.
\par\smallskip\textit{Formulation and concept references:} \href{https://arxiv.org/html/1910.07199v2}{[S1]}.

\subsection{Short English statement}
What is the exact reliable information rate when each transmitted bit independently disappears and the receiver sees only the surviving ordered subsequence?
\subsection{Sources}
\begin{itemize}
\item[S1]Cheraghchi and Ribeiro, An Overview of Capacity Results for Synchronization Channels.
\url{https://arxiv.org/html/1910.07199v2}.
\textit{Formulation source.}
\item[S2]Status source for Capacity of the Binary Deletion Channel.
\url{https://arxiv.org/abs/2609.19412}.
\textit{Further reference; not a proof endorsement.}
\end{itemize}
\end{document}
